\MulaiModul{07}{Array}

\subsection{Tujuan}

\begin{tujuan}
Setelah menyelesaikan modul ini, mahasiswa mampu:
\begin{enumerate}
    \item mendefinisikan data dengan \texttt{DB}, \texttt{DW}, dan
          \texttt{DUP};
    \item menjelaskan array sebagai data berurutan di memori dan
          string sebagai deretan karakter yang diakhiri \texttt{\$};
    \item menjelaskan tujuh mode pengalamatan yang dipakai pada modul
          ini;
    \item mengakses elemen array dengan \texttt{SI}, \texttt{BX},
          \texttt{[SI]}, dan \texttt{[BX]};
    \item menjumlahkan array word dan mencari elemen terbesar array
          byte;
    \item merakit dan menjalankan kedua program.
\end{enumerate}
\end{tujuan}

\subsection{Dasar singkat}

\subsubsection*{DB, DW, dan DUP}

\texttt{DB} mendefinisikan data satu byte, sedangkan \texttt{DW}
mendefinisikan data dua byte (satu word). \texttt{DUP} mengulang
sebuah nilai sebanyak yang ditentukan. Ketiga bentuk itu dirangkum
pada Tabel~\ref{tab:m07-data}.

\begin{table}[H]
    \centering
    \caption{Directif data yang dipakai pada Modul 07.}
    \label{tab:m07-data}
    \small
    \begin{tabular}{@{}llll@{}}
        \toprule
        \textbf{Directif} & \textbf{Ukuran} & \textbf{Contoh} &
        \textbf{Isi} \\
        \midrule
        \texttt{DB} & 1 byte & \texttt{DERET DB 5,9,3} &
        tiga byte berurutan \\
        \texttt{DW} & 2 byte & \texttt{DERET DW 2,3,4} &
        tiga word, total enam byte \\
        \texttt{DUP} & berulang & \texttt{RUANG DB 10 DUP(0)} &
        sepuluh byte bernilai nol \\
        \bottomrule
    \end{tabular}
\end{table}

Array adalah kumpulan data bertipe sama yang disimpan berurutan.
Elemen pertama berada pada offset awal array, elemen kedua pada
offset awal ditambah ukuran elemen, dan seterusnya. String adalah
deretan karakter yang setiap karakter menempati satu byte. String
yang akan dicetak dengan \texttt{AH=09h} diakhiri tanda \texttt{\$};
tanda itu menandai akhir string dan tidak ikut tercetak.

\subsubsection*{Tujuh mode pengalamatan}

Mode pengalamatan menyatakan bagaimana prosesor memperoleh operand.
Tujuh mode yang dipakai pada modul ini dirangkum pada
Tabel~\ref{tab:m07-mode}.

\begin{table}[H]
    \centering
    \caption{Tujuh mode pengalamatan pada Modul 07.}
    \label{tab:m07-mode}
    \small
    \begin{tabular}{@{}lll@{}}
        \toprule
        \textbf{Mode} & \textbf{Bentuk} & \textbf{Contoh} \\
        \midrule
        Register & register & \texttt{MOV AX,BX} \\
        Immediate & nilai segera & \texttt{MOV AX,5} \\
        Direct memory & nama variabel & \texttt{MOV AX,DERET} \\
        Indirect & \texttt{[register]} & \texttt{MOV AX,[SI]} \\
        Base relative & \texttt{[BX+disp]} & \texttt{MOV AX,[BX+2]} \\
        Direct index & \texttt{[SI+disp]} & \texttt{MOV AX,[SI+2]} \\
        Based indexed & \texttt{[BX+SI+disp]} &
        \texttt{MOV AX,[BX+SI+2]} \\
        \bottomrule
    \end{tabular}
\end{table}

Pada mode indirect dan seterusnya, \texttt{SI} atau \texttt{BX}
menyimpan alamat offset, bukan data. Tanda kurung siku memerintahkan
prosesor mengambil isi memori pada alamat yang ditunjuk register itu.
Penunjuk yang dipakai pada modul ini adalah \texttt{SI} dan
\texttt{BX}. Perpindahan ditulis sebagai bilangan biasa, dan besarnya
mengikuti ukuran elemen: dua byte untuk word, satu byte untuk byte.

\subsubsection*{Menelusuri array}

Pola penjumlahan array word memakai \texttt{SI} sebagai penunjuk dan
\texttt{CX} sebagai pencacah. Setiap iterasi menambahkan isi
\texttt{[SI]} ke akumulator, lalu \texttt{SI} ditambah 2 karena satu
elemen word berukuran dua byte. Jejak tiga iterasi ditunjukkan pada
Tabel~\ref{tab:m07-jejak}.

\begin{table}[H]
    \centering
    \caption{Jejak penjumlahan \texttt{DERET DW 2,3,4}.}
    \label{tab:m07-jejak}
    \small
    \begin{tabular}{@{}cccc@{}}
        \toprule
        \textbf{Iterasi} & \textbf{Elemen \texttt{[SI]}} &
        \textbf{AX sesudah} & \textbf{SI sesudah} \\
        \midrule
        1 & 2 & 2 & offset awal + 2 \\
        2 & 3 & 5 & offset awal + 4 \\
        3 & 4 & 9 & offset awal + 6 \\
        \bottomrule
    \end{tabular}
\end{table}

Pencarian nilai terbesar array byte memakai \texttt{BX} sebagai
penunjuk. Elemen pertama dijadikan kandidat pada \texttt{AL}. Setiap
elemen berikutnya dibandingkan dengan kandidat memakai
\texttt{CMP}; bila elemen lebih besar, kandidat diganti. Karena isi
array diperlakukan sebagai bilangan tidak bertanda, lompatan yang
dipakai adalah \texttt{JAE}. Setelah semua elemen diperiksa,
\texttt{AL} berisi nilai terbesar.

Jumlah yang melebihi 9 tidak dapat dicetak dengan menambahkan
\texttt{30h}. Nilai semacam itu diperiksa pada register dengan
emu8086 atau DEBUG. Modul ini hanya memakai array satu dimensi dan
tidak memakai \texttt{PUSH}, parameter lewat stack, atau pengurutan.

\begin{rujukan}
Rujukan singkat: Gunawan, Bagian 9 kecuali 9.7, dan Bagian 10.1;
S'to, Bab XI (mode yang dipakai pada modul ini); buku ajar,
\texttt{chapters/12} dan \texttt{chapters/15/section-01.tex};
buku ajar, \texttt{mode-array.asm} dan \texttt{array-string.asm}.
\end{rujukan}

\subsection{Alat dan bahan}

\begin{itemize}
    \item Komputer dengan emu8086 terpasang.
    \item DOSBox dengan TASM, TLINK, dan DEBUG.
    \item Berkas \texttt{code/m07-jumlah-array.asm} dan
          \texttt{code/m07-maksimum.asm}.
    \item Lembar jurnal praktikum.
\end{itemize}

\subsection{Contoh program}

Program pertama menjumlahkan tiga word dengan \texttt{SI} dan
\texttt{LOOP}. Jumlahnya 9 sehingga dapat dicetak sebagai satu digit.

\lstinputlisting[
    caption={Menjumlahkan array word dengan \texttt{SI} dan
    \texttt{LOOP}.},
    label={lst:m07-jumlah}
]{code/m07-jumlah-array.asm}

Program kedua mencari elemen terbesar array byte dengan \texttt{BX}.
Kandidat terbesar disimpan pada \texttt{AL} dan diperbarui memakai
\texttt{CMP} dan \texttt{JAE}.

\lstinputlisting[
    caption={Mencari elemen terbesar array byte dengan \texttt{BX}.},
    label={lst:m07-maksimum}
]{code/m07-maksimum.asm}

\subsection{Langkah kerja}

\begin{enumerate}
    \item Siapkan folder kerja, lalu salin kedua berkas \texttt{.asm}
          ke dalamnya.
    \item Rakit dan tautkan program pertama. Jalankan kedua perintah
          berikut pada baris terpisah:
          \begin{perintah}
tasm m07-jumlah-array.asm\\
tlink m07-jumlah-array.obj
\end{perintah}
          Jalankan hasilnya dan catat karakter yang tercetak.
    \item Jalankan program pertama dengan \emph{single step} pada
          emu8086. Catat \texttt{[SI]}, \texttt{AX}, dan \texttt{SI}
          pada setiap iterasi, lalu bandingkan dengan
          Tabel~\ref{tab:m07-jejak}.
    \item Ubah isi \texttt{DERET} menjadi \texttt{5,6,7} sehingga
          jumlahnya 18. Rakit ulang dan jalankan. Jelaskan mengapa
          keluarannya bukan bilangan 18, lalu periksa \texttt{AX}
          dengan emu8086 atau DEBUG.
    \item Rakit dan tautkan program kedua dengan perintah yang sama,
          lalu jalankan. Catat karakter yang tercetak.
    \item Jalankan program kedua dengan \emph{single step}. Catat
          \texttt{AL} dan \texttt{BX} pada setiap iterasi. Ubah isi
          \texttt{DERET} menjadi \texttt{4,8,2,6,1}, rakit ulang, dan
          pastikan keluarannya berubah.
\end{enumerate}

\subsection{Tugas}

\begin{enumerate}
    \item Jelaskan perbedaan \texttt{DB} dan \texttt{DW}, lalu
          tentukan jumlah byte untuk lima elemen masing-masing.
    \item Jelaskan tugas \texttt{DUP} dan berikan satu contoh beserta
          jumlah byte yang dihasilkan.
    \item Sebutkan tujuh mode pengalamatan pada modul ini dan berikan
          satu contoh untuk setiap mode.
    \item Jelaskan mengapa \texttt{SI} ditambah 2 pada array word,
          sedangkan \texttt{BX} ditambah 1 pada array byte.
    \item Jelaskan cara kerja pencarian nilai terbesar dan alasan
          pemakaian \texttt{JAE}.
\end{enumerate}

\subsection{Hasil yang dikumpulkan}

\begin{enumerate}
    \item Kedua berkas \texttt{.asm} beserta versi modifikasinya.
    \item Tangkapan layar keluaran kedua program.
    \item Tabel jejak iterasi berisi \texttt{[SI]}, \texttt{AX}, dan
          \texttt{SI} untuk program pertama serta \texttt{BX} dan
          \texttt{AL} untuk program kedua.
    \item Lembar jurnal berisi jawaban langkah kerja dan tugas.
\end{enumerate}
